% GENERATED FILE. Edit canonical lessons, then run npm run generate:books.
% canonical-source-hash: fnv1a64:af06ae3ef0836b61
% canonical-lessons: AR-C512-shifa, AR-C512-naqaha, AR-C512-taafa, AR-C512-afiya, AR-C512-taallama-pain

\chapter{Getting well again}
\label{ch:ar-getting-well-again}
\hlchaptermodality{512}

\begin{chapteropening}
\textbf{By the end of this chapter:} \emph{I can talk about recovery, convalescence and good health, and say that someone recovered or was in pain.}

Complete the last lesson of chapter 512: I can talk about recovery, convalescence and good health, and say that someone recovered or was in pain.
\end{chapteropening}

\section[\texorpdfstring{\ar{شفاء} (shifāʾ)}{شفاء (shifāʾ)}]{\ar{شفاء} (shifāʾ) — recovery, a cure}

\label{lesson:AR-C512-shifa}



\begin{quote}
Before the new one: say the Arabic for a splint, then the Arabic for to bind up.
\end{quote}

\subsection*{You'll want to know: \ar{شفاء}}
\textbf{\ar{شفاء}} — \emph{shifāʾ} — \textquotedblleft{}recovery, a cure\textquotedblright{}.

\begin{quote}
\ar{نتمنى} \ar{لك} \ar{الشفاء} \ar{العاجل}.
\end{quote}

\begin{quote}
\emph{natamannā laka ash-shifāʾ al-ʿājil.}
\end{quote}

\begin{quote}
We wish you a quick recovery.
\end{quote}

\ar{نتمنى} \ar{لك} \ar{الشفاء} is what you say to someone who is ill.

\begin{cousinweb}
\textbf{\ar{شفاء}} (\emph{shifāʾ}) is built on the native Arabic root \emph{sh-f-y}, \textquotedblleft{}to heal, to cure\textquotedblright{} (Lane). Shifāʾ is the healing itself.
\end{cousinweb}

\begin{grammarlens}[title={What you've built}]
One more word for describing what you lived and felt.
\end{grammarlens}

\subsection*{Your turn}
\begin{itemize}
\raggedright
  \item \emph{Say it:} \emph{shifāʾ}
  \item \emph{Say it:} \emph{shifāʾ}, once more
  \item \emph{Say it:} \emph{natamannā laka ash-shifāʾ al-ʿājil.} — We wish you a quick recovery.
  \item \emph{Recall:} say the Arabic for a splint, then the Arabic for to bind up, then say \emph{shifāʾ} again
\end{itemize}

\begin{tcolorbox}[breakable,colback=teal!4,colframe=teal!35!black,title={Before you move on}]
What does \ar{شفاء} mean? (\textquotedblleft{}Recovery, a cure\textquotedblright{}.) Say it once more.
\end{tcolorbox}
\section[\texorpdfstring{\ar{نقاهة} (naqāha)}{نقاهة (naqāha)}]{\ar{نقاهة} (naqāha) — convalescence}

\label{lesson:AR-C512-naqaha}



\begin{quote}
Before the new one: say the Arabic for to bind up, then the Arabic for recovery.
\end{quote}

\subsection*{You'll want to know: \ar{نقاهة}}
\textbf{\ar{نقاهة}} — \emph{naqāha} — \textquotedblleft{}convalescence\textquotedblright{}.

\begin{quote}
\ar{بقي} \ar{أبي} \ar{في} \ar{البيت} \ar{فترة} \ar{النقاهة}.
\end{quote}

\begin{quote}
\emph{baqiya abī fī al-bayt fatrat an-naqāha.}
\end{quote}

\begin{quote}
My father stayed at home during his convalescence.
\end{quote}

\begin{cousinweb}
\textbf{\ar{نقاهة}} (\emph{naqāha}) is usually traced to the native Arabic root \emph{n-q-h}, \textquotedblleft{}to recover from an illness\textquotedblright{}: the weak days between sickness and full health.
\end{cousinweb}

\begin{grammarlens}[title={What you've built}]
One more word for describing what you lived and felt.
\end{grammarlens}

\subsection*{Your turn}
\begin{itemize}
\raggedright
  \item \emph{Say it:} \emph{naqāha}
  \item \emph{Say it:} \emph{naqāha}, once more
  \item \emph{Say it:} \emph{baqiya abī fī al-bayt fatrat an-naqāha.} — My father stayed at home during his convalescence.
  \item \emph{Recall:} say the Arabic for to bind up, then the Arabic for recovery, then say \emph{naqāha} again
\end{itemize}

\begin{tcolorbox}[breakable,colback=teal!4,colframe=teal!35!black,title={Before you move on}]
What does \ar{نقاهة} mean? (\textquotedblleft{}Convalescence\textquotedblright{}.) Say it once more.
\end{tcolorbox}
\section[\texorpdfstring{\ar{تعافى} (taʿāfā)}{تعافى (taʿāfā)}]{\ar{تعافى} (taʿāfā) — to recover}

\label{lesson:AR-C512-taafa}



\begin{quote}
Before the new one: say the Arabic for recovery, then the Arabic for convalescence.
\end{quote}

\subsection*{You'll want to know: \ar{تعافى}}
\textbf{\ar{تعافى}} — \emph{taʿāfā} — \textquotedblleft{}to recover\textquotedblright{}.

\begin{quote}
\ar{تعافى} \ar{اللاعب} \ar{وعاد} \ar{إلى} \ar{الفريق}.
\end{quote}

\begin{quote}
\emph{taʿāfā al-lāʿib wa-ʿāda ilā al-farīq.}
\end{quote}

\begin{quote}
The player recovered and went back to the team.
\end{quote}

Present: \ar{يتعافى}, \emph{yataʿāfā}.

\begin{cousinweb}
\textbf{\ar{تعافى}} (\emph{taʿāfā}) is built on the native Arabic root \emph{ʿ-f-w}, \textquotedblleft{}to efface; to pardon; to give health\textquotedblright{} (Lane). ʿAfwan, which you know, is from the same root.
\end{cousinweb}

\begin{grammarlens}[title={What you've built}]
One more word for describing what you lived and felt.
\end{grammarlens}

\subsection*{Your turn}
\begin{itemize}
\raggedright
  \item \emph{Say it:} \emph{taʿāfā}
  \item \emph{Say it:} \emph{taʿāfā}, once more
  \item \emph{Say it:} \emph{taʿāfā al-lāʿib wa-ʿāda ilā al-farīq.} — The player recovered and went back to the team.
  \item \emph{Recall:} say the Arabic for recovery, then the Arabic for convalescence, then say \emph{taʿāfā} again
\end{itemize}

\begin{tcolorbox}[breakable,colback=teal!4,colframe=teal!35!black,title={Before you move on}]
What does \ar{تعافى} mean? (\textquotedblleft{}To recover\textquotedblright{}.) Say it once more.
\end{tcolorbox}
\section[\texorpdfstring{\ar{عافية} (ʿāfiya)}{عافية (ʿāfiya)}]{\ar{عافية} (ʿāfiya) — good health, wellbeing}

\label{lesson:AR-C512-afiya}



\begin{quote}
Before the new one: say the Arabic for convalescence, then the Arabic for to recover.
\end{quote}

\subsection*{You'll want to know: \ar{عافية}}
\textbf{\ar{عافية}} — \emph{ʿāfiya} — \textquotedblleft{}good health, wellbeing\textquotedblright{}.

\begin{quote}
\ar{الحمد} \ar{لله،} \ar{جدي} \ar{في} \ar{عافية}.
\end{quote}

\begin{quote}
\emph{al-ḥamdu lillāh, jaddī fī ʿāfiya.}
\end{quote}

\begin{quote}
Thank God, my grandfather is in good health.
\end{quote}

\begin{cousinweb}
\textbf{\ar{عافية}} (\emph{ʿāfiya}) is built on the native Arabic root \emph{ʿ-f-w}, \textquotedblleft{}to efface; to pardon; to give health\textquotedblright{} (Lane). ʿĀfiya is health given back, freedom from illness.
\end{cousinweb}

\begin{grammarlens}[title={What you've built}]
One more word for describing what you lived and felt.
\end{grammarlens}

\subsection*{Your turn}
\begin{itemize}
\raggedright
  \item \emph{Say it:} \emph{ʿāfiya}
  \item \emph{Say it:} \emph{ʿāfiya}, once more
  \item \emph{Say it:} \emph{al-ḥamdu lillāh, jaddī fī ʿāfiya.} — Thank God, my grandfather is in good health.
  \item \emph{Recall:} say the Arabic for convalescence, then the Arabic for to recover, then say \emph{ʿāfiya} again
\end{itemize}

\begin{tcolorbox}[breakable,colback=teal!4,colframe=teal!35!black,title={Before you move on}]
What does \ar{عافية} mean? (\textquotedblleft{}Good health, wellbeing\textquotedblright{}.) Say it once more.
\end{tcolorbox}
\section[\texorpdfstring{\ar{تألم} (taʾallama)}{تألم (taʾallama)}]{\ar{تألم} (taʾallama) — to be in pain}

\label{lesson:AR-C512-taallama-pain}



\begin{quote}
Before the new one: say the Arabic for to recover, then the Arabic for good health.
\end{quote}

\subsection*{You'll want to know: \ar{تألم}}
\textbf{\ar{تألم}} — \emph{taʾallama} — \textquotedblleft{}to be in pain\textquotedblright{}.

\begin{quote}
\ar{تألم} \ar{الرجل} \ar{كثيرا} \ar{من} \ar{ظهره}.
\end{quote}

\begin{quote}
\emph{taʾallama ar-rajul kathīran min ẓahrihi.}
\end{quote}

\begin{quote}
The man suffered a lot from his back.
\end{quote}

Present: \ar{يتألم}, \emph{yataʾallam}.

It takes \ar{من} for the place: \ar{تألم} \ar{من} \ar{ظهره}.

\begin{cousinweb}
\textbf{\ar{تألم}} (\emph{taʾallama}) is built on the native Arabic root \emph{ʾ-l-m}, \textquotedblleft{}pain\textquotedblright{} (Lane). Alam, pain, which you know, is from the same root; Form V is to feel it.
\end{cousinweb}

\begin{grammarlens}[title={What you've built}]
One more word for describing what you lived and felt.
\end{grammarlens}

\subsection*{Your turn}
\begin{itemize}
\raggedright
  \item \emph{Say it:} \emph{taʾallama}
  \item \emph{Say it:} \emph{taʾallama}, once more
  \item \emph{Say it:} \emph{taʾallama ar-rajul kathīran min ẓahrihi.} — The man suffered a lot from his back.
  \item \emph{Recall:} say the Arabic for to recover, then the Arabic for good health, then say \emph{taʾallama} again
\end{itemize}

\begin{tcolorbox}[breakable,colback=teal!4,colframe=teal!35!black,title={Before you move on}]
What does \ar{تألم} mean? (\textquotedblleft{}To be in pain\textquotedblright{}.) Say it once more.
\end{tcolorbox}
